\scanpage{19}\mspage{4}
utilise (1.1), en regardant les deux membres de (1.3) comme des sections de $\mathrm{p}_{*}\handwritten{\mathcal{O}}(6\mathrm{e})$ : en effet, $\mathrm{y}^{2}-\mathrm{x}^{3}$ est dans $\mathrm{p}_{*}\handwritten{\mathcal{O}}(5\mathrm{e})$, $\mathrm{xy}$ dans $\mathrm{p}_{*}\handwritten{\mathcal{O}}(5\mathrm{e})$, $\mathrm{x}^{2}$ dans $\mathrm{p}_{*}\handwritten{\mathcal{O}}(4\mathrm{e})$, $\mathrm{y}$ dans $\mathrm{p}_{*}\handwritten{\mathcal{O}}(3\mathrm{e})$, $\mathrm{x}$ dans $\mathrm{p}_{*}\handwritten{\mathcal{O}}(2\mathrm{e})$ et $1$ dans $\mathrm{p}_{*}\handwritten{\mathcal{O}}$.

Si on pose
\[
\mathrm{F}(\mathrm{x},\mathrm{y})=\mathrm{y}^{2}+\mathrm{a}_{1}\,\mathrm{xy}+\mathrm{a}_{3}\,\mathrm{y}-\mathrm{x}^{3}-\mathrm{a}_{2}\,\mathrm{x}^{2}-\mathrm{a}_{4}\,\mathrm{x}-\mathrm{a}_{6}
\]
la forme différentielle invariante définie par $\pi$ s'écrit
\[
\begin{aligned}
\pi&=\mathrm{dx}/\mathrm{F}_{\mathrm{y}}=-\mathrm{dy}/\mathrm{F}_{\mathrm{x}}\qquad\text{soit}\\[.65em]
\pi&=\handwritten{\frac}{\mathrm{dx}}{2\mathrm{y}+\mathrm{a}_{1}\mathrm{x}+\mathrm{a}_{3}}
=\handwritten{\frac}{\mathrm{dy}}{3\mathrm{x}^{2}+2\mathrm{a}_{2}\mathrm{x}+\mathrm{a}_{4}-\mathrm{a}_{1}\mathrm{y}}\qquad.
\end{aligned}
\]

Réciproquement toute équation (1.3) est l'équation non homogène d'une courbe de genre $1$ plongée comme une cubique dans $\mathrm{P}_{2}$ et présentant un point d'inflexion en $(0,1,0)$\margin{\handwritten{$\Big)$\,?}} avec la droite de l'infini pour tangente d'inflexion.

On pose alors
% A short isolated stroke at left faces the row beginning d_6 below.
\reversemarginpar
\[
\mathrm{b}_{2}=\mathrm{a}_{1}^{2}+4\mathrm{a}_{2}\qquad
\mathrm{b}_{4}=\mathrm{a}_{1}\mathrm{a}_{3}+2\mathrm{a}_{4}\qquad
\mathrm{b}_{6}=\mathrm{a}_{3}^{2}+4\mathrm{a}_{6}
\tag*{(1.4)}
\]
\[
\begin{aligned}
&-\mathrm{b}_{8}=\mathrm{a}_{1}\,\mathrm{a}_{3}\,\mathrm{a}_{4}+\mathrm{a}_{4}^{2}-\mathrm{a}_{1}^{2}\,\mathrm{a}_{6}-\mathrm{a}_{2}\,\mathrm{a}_{3}^{2}-\overset{\inserted{\handwritten{2^{2}}}}{4}\mathrm{a}_{2}\,\mathrm{a}_{6}\\[.55em]
&\mathrm{d}_{6}=\mathrm{b}_{2}\,\mathrm{b}_{4}-\overset{\inserted{\handwritten{2^{1}.3^{2}}}}{18}\mathrm{b}_{6}\qquad
\mathrm{d}_{8}=\mathrm{b}_{2}^{2}\,\mathrm{b}_{4}-2^{4}.3\,\mathrm{b}_{4}^{2}+2.3^{2}\,\mathrm{b}_{2}\,\mathrm{b}_{6}
\end{aligned}
\]
\margin{\handwritten{$\backslash$}}\normalmarginpar

On a
\[
\begin{aligned}
&-4\mathrm{b}_{8}=\mathrm{b}_{4}^{2}-\mathrm{b}_{2}\,\mathrm{b}_{6}\\[.45em]
&\tikz[baseline=(digit.base)]{\node[inner sep=0pt] (digit) {$1$};\handwritten{\draw[thin,gray,rotate around={-12:(digit.center)}] (digit.center) ellipse [x radius=.43em,y radius=1.32em];}}2\,\mathrm{d}_{6}=\mathrm{c}_{6}+\mathrm{b}_{2}\,\mathrm{c}_{4}\qquad
-\tikz[baseline=(digit.base)]{\node[inner sep=0pt] (digit) {$1$};\handwritten{\draw[thin,gray,rotate around={-12:(digit.center)}] (digit.center) ellipse [x radius=.4em,y radius=1.13em];}}2\,\mathrm{d}_{8}=\mathrm{c}_{4}^{2}+\mathrm{b}_{2}\,\mathrm{c}_{6}
\end{aligned}
\]
\clearpage
% archive sheet number: 18
